\documentclass[12pt,a4paper]{article}
\usepackage[utf8]{inputenc}
\usepackage[T2A]{fontenc}
\usepackage[russian]{babel}
\usepackage{amsmath,amssymb}
\usepackage{booktabs,array}
\usepackage{geometry}
\usepackage{microtype}
\geometry{margin=2.3cm}

\title{Continuous Wilson gap-action gate\\
устойчивость угла и геометрическое сокращение осей}
\author{}
\date{5 августа 2026 г.}

\begin{document}
\maketitle

\section{Непрерывное двухсекторное решение}

Для \(SO(3)\) Wilson-holonomy с фазами \((0,a,1-a)\) полный
tensor-response в единицах \(\pi^4\) равен
\begin{equation}
 R(c)=\frac1{45}+\frac{8(2+c)}{3(1-c)^2},
 \qquad
 c=\cos(2\pi a).
\end{equation}
Условие \(R(c)=1\) даёт
\begin{equation}
 11c^2-52c-49=0,
\end{equation}
и допустимый корень
\begin{equation}
 c_*=\frac{26-9\sqrt{15}}{11},
 \qquad
 a_*\approx0.3989625047.
\end{equation}
Перебор канонических octahedral-\(S_4\) осей даёт восемь направлений,
которые одновременно сохраняют \(R(c_*)=1\) и расширяют family-алгебру до
\(M_3\).

\section{Spectral gap-functional}

Точка \(c_*\) является не только решением обратного численного условия.
Функция
\begin{equation}
 V_{\mathrm{gap}}(c)
 =
 \frac83\left(
 \frac3{1-c}+\log(1-c)
 \right)
 -\frac{44}{45}c
\end{equation}
удовлетворяет
\begin{equation}
 V_{\mathrm{gap}}'(c)=R(c)-1.
\end{equation}
Поэтому \(c_*\) является её стационарной точкой. Более того,
\begin{equation}
 V_{\mathrm{gap}}''(c_*)
 =R'(c_*)
 =\frac{8(5+c_*)}{3(1-c_*)^3}
 \approx1.901648596>0.
\end{equation}
Следовательно, найденный угол является устойчивым локальным минимумом
явного zero-parameter spectral primitive.

Это повышает статус конструкции: угол можно читать как решение gap-equation,
а не как свободно назначенный holonomy. Но единица в условии
\(R(c)=1\) теперь становится новым normalization gate. При замене единицы на
bare stiffness \(\kappa\)
\begin{equation}
 \frac{dc_*}{d\kappa}\bigg|_{\kappa=1}
 =\frac1{R'(c_*)}
 \approx0.525859511.
\end{equation}
Поэтому каноническая нормировка должна быть выведена, а не объявлена.

\section{Factor-axis selector}

К оси \(\hat n\) применены уже существующие factor operators
\begin{equation}
 L(w_3,w_1)
 =w_3(I-T_{\mathbb{RP}^3})+w_1(I-T_{S^1})
\end{equation}
с тремя ранее объявленными kernels: inverse length, inverse square и
tunneling. Для каждого из них минимизация
\begin{equation}
 E_L(\hat n)=\hat n^TL\hat n
\end{equation}
сокращает восемь совместных осей до двух эквивалентных
transposition-направлений. Таким образом, старая product geometry уже
содержит частичный selector:
\begin{equation}
 8\longrightarrow2.
\end{equation}

Оставшаяся двукратность может интерпретироваться как symmetry-related vacuum
degeneracy, но это допустимо только после вывода самого coupling
\(\hat n^TL\hat n\) из того же boundary action.

\section{Вердикт}

\begin{equation}
 \boxed{
 \begin{gathered}
 \text{joint operator existence}\quad\checkmark,\\
 \text{stable spectral saddle}\quad\checkmark,\\
 \text{partial geometric axis selector}\quad\checkmark,\\
 \text{единый local/BV action и unit stiffness}\quad\text{открыты}.
 \end{gathered}
 }
\end{equation}

Следующий строгий шаг: построить auxiliary-field или boundary gauge action,
чьё однопетлевое интегрирование одновременно даёт
\(V_{\mathrm{gap}}\), каноническую единичную stiffness и factor-axis term.

Машинный аудит сохранён в
\texttt{s2t\_continuous\_wilson\_gap\_action\_audit.py} и
\texttt{s2t\_continuous\_wilson\_gap\_action\_results.json}.

\end{document}